\section{Résultats}
\label{sec:results}

\subsection{Familles de caractéristiques individuelles}

La Figure~\ref{fig:individual} présente les résultats du benchmark pour chaque
jeu mono-famille. Parmi les quatre, \textbf{D --- HOG + SVM} obtient la meilleure
performance individuelle (\AccDHog{}, F1-macro \FOneDHog), suivi de B --- LBP (\AccBTex),
A --- Stats (\AccAStats) et C --- DWT (\AccCDwt).

\begin{figure}[H]
  \centering
  \begin{subfigure}[b]{0.49\linewidth}
    \includegraphics[width=\linewidth]{p2_A_—_Stats.png}
    \caption{A --- Statistiques (meilleur : \ModelAStats, \AccAStats)}
  \end{subfigure}
  \hfill
  \begin{subfigure}[b]{0.49\linewidth}
    \includegraphics[width=\linewidth]{p2_B_—_Tex(opt).png}
    \caption{B --- Texture LBP (meilleur : \ModelBTex, \AccBTex)}
  \end{subfigure}
  \begin{subfigure}[b]{0.49\linewidth}
    \includegraphics[width=\linewidth]{p2_C_—_DWT(opt).png}
    \caption{C --- Fréquence DWT (meilleur : \ModelCDwt, \AccCDwt)}
  \end{subfigure}
  \hfill
  \begin{subfigure}[b]{0.49\linewidth}
    \includegraphics[width=\linewidth]{p2_D_—_HOG.png}
    \caption{D --- HOG (meilleur : \ModelDHog, \AccDHog)}
  \end{subfigure}
  \caption{Résultats du benchmark Phase~2 pour les familles de caractéristiques individuelles.}
  \label{fig:individual}
\end{figure}

\subsection{Jeux de caractéristiques fusionnés}

La fusion progressive des caractéristiques améliore systématiquement les performances
(Figure~\ref{fig:combined}). Le gain est le plus marqué lorsque le HOG est inclus
dans le jeu complet.

\begin{figure}[H]
  \centering
  \begin{subfigure}[b]{0.49\linewidth}
    \includegraphics[width=\linewidth]{p2_ApB.png}
    \caption{A+B : Statistiques + LBP}
  \end{subfigure}
  \hfill
  \begin{subfigure}[b]{0.49\linewidth}
    \includegraphics[width=\linewidth]{p2_BpC(opt).png}
    \caption{B+C : LBP + DWT}
  \end{subfigure}
  \begin{subfigure}[b]{0.49\linewidth}
    \includegraphics[width=\linewidth]{p2_ApBpC(opt).png}
    \caption{A+B+C : sans HOG}
  \end{subfigure}
  \hfill
  \begin{subfigure}[b]{0.49\linewidth}
    \includegraphics[width=\linewidth]{p2_Full(opt).png}
    \caption{Full(opt) : les quatre familles}
  \end{subfigure}
  \caption{Résultats du benchmark Phase~2 pour les jeux de caractéristiques combinés.}
  \label{fig:combined}
\end{figure}

\subsection{Classement global}

Le Tableau~\ref{tab:leaderboard} rapporte les cinq meilleures configurations.
Le diagramme en barres complet est présenté à la Figure~\ref{fig:leaderboard}.

\begin{table}[H]
  \centering
  \caption{Top 5 des configurations classées par F1-macro sur le test tenu à l'écart.}
  \label{tab:leaderboard}
  \begin{tabular}{@{}clcccc@{}}
    \toprule
    \textbf{Rang} & \textbf{Jeu de caract.} & \textbf{Classifieur} & \textbf{F1-macro} & \textbf{Précision} & \textbf{$\kappa$} \\
    \midrule
    1 & Full(opt)   & SVM       & \textbf{\BestFOne} & \textbf{\BestAcc} & \textbf{\BestKappa} \\
    2 & D --- HOG   & SVM       & \RankTwoFOne  & \HogSvmAcc  & 0.855 \\
    3 & Full(opt)   & KNN       & \RankThreeFOne & 90.0\%      & 0.850 \\
    4 & D --- HOG   & KNN       & \RankFourFOne  & 89.7\%      & 0.845 \\
    5 & Full(opt)   & LightGBM  & \RankFiveFOne  & 89.3\%      & 0.840 \\
    \bottomrule
  \end{tabular}
\end{table}

\begin{figure}[H]
  \centering
  \includegraphics[width=0.88\linewidth]{p2_leaderboard.png}
  \caption{Diagramme en barres du classement de toutes les configurations évaluées,
           triées par F1-macro. Full(opt) + SVM est le pipeline le mieux classé.}
  \label{fig:leaderboard}
\end{figure}

Le tableau de bord global (Figure~\ref{fig:dashboard}) synthétise le F1-macro,
le temps d'entraînement et la variabilité par jeu de caractéristiques pour tous
les classifieurs.

\begin{figure}[H]
  \centering
  \includegraphics[width=\linewidth]{p2_overall_dashboard.png}
  \caption{Tableau de bord Phase~2 : F1-macro et temps d'entraînement pour tous
           les jeux de caractéristiques et classifieurs.}
  \label{fig:dashboard}
\end{figure}
